\section{鲁棒与准确不可兼得}
\label{sec:17-Constrained-Guarantees/02-Adversarial-Robustness/03}

一个团队用对抗训练把模型的鲁棒准确率从 \(4\%\) 提到 \(52\%\)，干净数据上的准确率从 \(96\%\) 掉到 \(87\%\)。产品那边不接受掉九个点，于是开始调：换结构、加数据、调学习率、试各种正则、混合干净样本与对抗样本的比例。三个月过去，最好的一版是鲁棒 \(53\%\)、干净 \(88\%\)。

这三个月不算白费，但它们回答的是一个已经有答案的问题。干净准确率与鲁棒准确率不能同时拉满，这件事有构造性反例，属于结论而不是现象。

\subsection{一个能把话说死的构造}

构造要简单到能算清楚，同时保留真实数据的关键特征：既有强但脆弱的信号，也有弱但可靠的信号。

设标签 \(y\in\{-1,+1\}\) 均匀取值。输入 \(x\in\R^{d+1}\) 由两部分组成。

第一个分量是一个\textbf{可靠但不完美}的特征：\(x_1=y\) 的概率为 \(p\)，\(x_1=-y\) 的概率为 \(1-p\)。取 \(p=0.95\)。它与标签强相关，但天花板是 \(95\%\)。

其余 \(d\) 个分量是一批\textbf{各自很弱、合起来很强}的特征：\(x_2,\dots,x_{d+1}\) 独立同分布于 \(\mathcal N(\eta y,\,1)\)，其中 \(\eta=4/\sqrt d\)。单看任何一个分量，信号 \(\eta\) 远小于噪声 1，几乎没用；把 \(d\) 个平均起来，均值服从 \(\mathcal N(\eta y,\,1/d)\)，信噪比是 \(\eta\sqrt d=4\)，判对的概率是 \(\Phi(4)\approx0.99997\)。

于是有两类分类器。只看 \(x_1\) 的，准确率 \(95\%\)。看后 \(d\) 个分量平均值的，准确率 \(99.997\%\)。标准训练当然会选后者——它的准确率高得多，而训练目标里没有任何一项反对这个选择。

现在让对手在 \(\ell_\infty\) 意义下扰动，预算 \(\varepsilon=2\eta\)。这个预算很小：每个分量只能挪动 \(8/\sqrt d\)，\(d\) 取几千时是零点几，远小于噪声的标准差 1，肉眼在数据里根本看不出来。但它足够做一件事——把后 \(d\) 个分量整体平移 \(-2\eta y\)。平移之后，它们的分布从 \(\mathcal N(\eta y,1)\) 变成 \(\mathcal N(-\eta y,1)\)，也就是\textbf{恰好等于相反类别的分布}。那个 \(99.997\%\) 的分类器会以同样的把握给出相反的答案。而 \(x_1\) 是离散的、只取 \(\pm1\)，这个扰动碰不到它。

\begin{theorem}[标准与鲁棒的和被卡住]
在上述分布上，取 \(\varepsilon=2\eta\)。任何分类器 \(f\)（确定性或随机的都可以）满足
\[
\mathrm{Acc}_{\mathrm{std}}(f)+\mathrm{Acc}_{\mathrm{rob}}(f)\;\le\;2p .
\]
\(p=0.95\) 时右端是 \(1.90\)：两个准确率不可能同时超过 \(95\%\)。
\end{theorem}

证明骨架。只用上面那一个特定的对手（整体平移后 \(d\) 个分量），得到的鲁棒准确率是真实鲁棒准确率的上界，所以证这个对手下的不等式就够。关键一步是前面那句话的形式化：给定 \(y\) 与 \(x_1\)，平移后的后 \(d\) 个分量与\textbf{类别 \(-y\) 的干净样本}同分布。于是把分类器的行为整理成四个量——记 \(q(a,m)=P\big(f\text{ 输出}+1\;\big|\;x_1=a,\ \text{后 }d\text{ 个分量来自类别 }m\big)\)，标准准确率与鲁棒准确率都是这四个量的线性组合，二者相加时含 \(q\) 的项按 \(x_1\) 的条件概率 \(p\) 与 \(1-p\) 配对相消，只剩
\[
\mathrm{Acc}_{\mathrm{std}}+\mathrm{Acc}_{\mathrm{rob}}
=1+\frac{2p-1}{2}\Big(q(1,1)+q(1,-1)-q(-1,1)-q(-1,-1)\Big).
\]
括号里的量至多是 \(2\)（前两项各不超过 1，后两项各不小于 0），代入即得 \(1+(2p-1)=2p\)。

条件对账值得逐条看，因为它决定了这个构造能推广到什么程度。\(2p\) 这个上界里 \(p\) 是那个抗扰动特征的可靠度上限，构造里它被硬编码成 \(0.95\)；换一个分布，卡住的位置随之改变，但只要存在``强而脆弱''与``弱而可靠''的分工，就会有一个小于 2 的上界。\(\varepsilon=2\eta\) 这个取值是让脆弱特征恰好被翻转所需的量，\(\varepsilon\) 更小时不等式会松弛，构造给出的限制随之减弱——所以这不是一条对所有 \(\varepsilon\) 都成立的普适定理，它说的是``存在分布和存在扰动预算，使得两者不可兼得''。这已经足够推翻``只要调参得当就能两全''这个默认前提，这正是反例的用途。

\begin{remark}
上界取到等号的分类器是只看 \(x_1\) 的那一个：它的标准准确率与鲁棒准确率都是 \(p=0.95\)，和恰好是 \(2p\)。而标准训练选出的那个分类器落在另一个极端，标准准确率 \(0.99997\)、鲁棒准确率约 \(0.00003\)，和恰好是 1。两者都在同一条约束线的允许范围内，区别在于工作点选在哪里——这是选择，不是优劣。
\end{remark}

\begin{center}
\begin{tikzpicture}[x=0.62cm,y=3.2cm,>=stealth,font=\small]
  \fill[green!12] (4,0) rectangle (8,1.02);
  \draw[->,black!60] (-0.5,0) -- (17.4,0)
    node[right,font=\footnotesize] {\(\varepsilon\)};
  \draw[->,black!60] (0,-0.03) -- (0,1.12)
    node[above,font=\footnotesize] {准确率};
  \foreach \e in {4,8,12,16} {
    \draw[black!60] (\e,0) -- (\e,-0.02);
    \node[below,font=\footnotesize] at (\e,-0.025) {\(\e/255\)};
  }
  \foreach \a/\lab in {0.5/{0.5}, 1.0/{1.0}} {
    \draw[black!60] (0,\a) -- (-0.35,\a);
    \node[left,font=\footnotesize] at (-0.4,\a) {\lab};
  }
  \draw[thick,blue!65] plot[domain=0:16.5,samples=40] (\x,{0.94-0.014*\x});
  \draw[thick,green!45!black] plot[domain=0:16.5,samples=60]
    (\x,{0.88*exp(-0.054*\x)});
  \draw[thick,orange!85!black] plot[domain=0:6,samples=70]
    (\x,{0.96*exp(-1.35*\x)});
  \draw[thick,orange!85!black] (6,0.002) -- (16.5,0.002);
  \node[anchor=west,font=\footnotesize,blue!65] at (8.4,0.90)
    {对抗训练：干净准确率};
  \node[anchor=west,font=\footnotesize,green!45!black] at (9.9,0.26)
    {对抗训练：鲁棒准确率};
  \node[anchor=west,font=\footnotesize,orange!85!black] at (2.3,0.28)
    {标准训练：鲁棒准确率};
  \node[font=\footnotesize,black!70] at (6,1.075) {常用工作区间};
  \draw[<->,black!70] (8,0.578) -- (8,0.828);
  \node[anchor=west,font=\footnotesize,black!70] at (8.2,0.70)
    {同一模型上的两个数};
\end{tikzpicture}
\end{center}

\emph{要记住的是橙线的形状：标准训练的模型不是``鲁棒性差一点''，是在很小的扰动预算下就归零。对抗训练把它抬起来，抬起来的同时蓝线在往下走——那个双向箭头的长度不会随着调参消失，它是上面那条不等式在具体模型上的样子。}

\subsection{对抗训练在优化一个自己算不准的目标}

既然要在这条线上选点，就得有办法选。对抗训练把经验风险最小化换成一个极小极大问题：
\[
\min_\theta\ \E_{(x,y)}\Big[\max_{\norm{\delta}\le\varepsilon}\ \ell\big(f_\theta(x+\delta),\,y\big)\Big].
\]
外层照旧是对参数最小化，内层是对扰动最大化。写成这个形状之后，``模型在邻域内都要对''这件事第一次进入了训练目标——上一节说过，训练目标没要求的事情优化器不会顺带做到，反过来，要它做到就得写进目标。

问题在内层。\(\ell\circ f_\theta\) 关于 \(\delta\) 一般非凸，那个 \(\max\) 没有闭式解，也没有可靠的求解算法。实际做法是从随机起点出发做几步投影梯度上升，得到一个 \(\hat\delta\)，用 \(\ell(f_\theta(x+\hat\delta),y)\) 代替真正的最大值。

这一步的性质要说清楚：\(\hat\delta\) 是可行解，所以
\[
\ell\big(f_\theta(x+\hat\delta),y\big)\;\le\;\max_{\norm{\delta}\le\varepsilon}\ell\big(f_\theta(x+\delta),y\big).
\]
也就是说，训练时实际最小化的是真实鲁棒风险的一个\textbf{下界}。最小化下界不等于最小化原目标——下界可以被压到很低，而原目标纹丝不动，只要那几步梯度上升恰好找不到真正的最坏扰动。这个逻辑与\hyperref[sec:09-Convex-Optimization/05-Duality/01]{``Lagrangian 怎样制造全局下界''}里的结构相反而同源：那里下界是用来证明最优性的证书，这里下界是被拿来当目标本身用的，两者能不能等同取决于间隙闭不闭合，而\hyperref[sec:09-Convex-Optimization/05-Duality/02]{``对偶间隙何时闭合''}给出的条件在这里一个都不满足。

同样的形状在\hyperref[sec:14-Deep-Learning/07-Generative-Models-Unified/03]{``GAN 把散度的估计交给判别器''}里出现过：内层的最优化被一个训练不足的近似替代，外层于是在优化一个自己不知道有多准的量。两处的现象也相似——内层算得不够狠时，外层会得到虚高的表现。

由此可以理解几条实践中反复出现的经验规律。攻击步数太少时训练出的模型往往在强攻击下失守，因为它只学会了躲开弱攻击找得到的那些扰动；对抗训练需要的模型容量和数据量明显高于标准训练；训练不稳定的现象比标准训练常见，因为目标函数每一步都在变（内层的解随参数移动）。这些都是经验规律，没有对应的收敛定理；而\hyperref[sec:09-Convex-Optimization/08-Nonconvex-and-Stochastic-Optimization/01]{``下降引理在非凸地形中只保证接近驻点''}那条判断在这里依然适用，只是连``目标函数''本身都只有近似值可用。

\subsection{梯度遮蔽：看起来鲁棒，实际不是}

对抗训练很贵，于是不断有更便宜的方案出现，其中一大类的共同特征是：让梯度不可用。

具体手法五花八门。在输入端加一个不可微的预处理（量化、JPEG 压缩、随机裁剪），基于梯度的攻击拿不到有意义的导数。在网络里插入随机性，每次前向的梯度都不一样，攻击的更新方向互相抵消。用极陡的激活让梯度爆炸或消失，数值上算出来的梯度是噪声。

这些方案在评测里的表现常常很漂亮：跑标准的梯度攻击，准确率几乎不掉。然后换一种攻击就全部失效——用可微的近似替代那个不可微的模块，或者对随机性取期望之后再求梯度，或者干脆用不需要梯度的黑盒攻击、用另一个模型上构造的扰动做迁移。这类``防御''在过去几年被系统性地推翻过好几轮，几乎没有例外。

问题的根源不在具体手法，在评判标准。``我们跑了攻击，没打穿''这句话的逻辑地位，和\hyperref[sec:11-Mathematical-Statistics/05-Hypothesis-Testing/01]{``两类错误与功效共同定义检验''}里那条提醒完全一样：没有拒绝原假设不等于原假设为真，尤其当检验的功效本来就很低时。梯度遮蔽做的事情，正是把攻击的功效降到接近零——它没有让模型更安全，只是让这个特定的检验失去了发现问题的能力。

\begin{proposition}[判断一个防御是否真实]
把\hyperref[sec:01-Mathematical-Language/01-Logic/03]{``量词与否定''}里的形状套上来：鲁棒性是``对邻域内所有扰动都正确''，它的否定是``存在一个扰动使预测出错''。攻击是在找存在性证据，找到了就推翻了鲁棒性；找不到，什么也没证明——全称命题不能靠有限次搜索确立。能确立它的只有第二节那种论证：给出一个上界（Lipschitz 常数、或高斯平滑后的可控性），由它推出邻域内不可能变号。
\end{proposition}

所以判据很简单：这个防御给出的是一个可证明的界，还是一份攻击失败的记录？前者哪怕数字难看也是承诺，后者哪怕数字漂亮也只是观察。二者不能互换，也不能互相佐证。

这与全书的``条件对账''是同一件事。一个结论的效力来自它依赖的条件被逐条核对过，而不是来自它在若干次尝试中没有被推翻。\hyperref[sec:01-Mathematical-Language/03-Proof-Techniques/04]{``怎样发现并修复错误证明''}里那套习惯——找出每一步在用哪个条件，然后问这个条件失效时会发生什么——在这里可以直接照搬：梯度遮蔽类防御的``证明''里，那一步隐含的条件是``攻击者只会用梯度''，而这个条件从来没被写下来，也从来不成立。

\subsection{这个单元留下的判断}

三节走下来，对抗鲁棒性的形状清楚了。

脆弱性的来源是几何，不是缺陷。高维空间里，一个 \(\ell_\infty\) 球的对角伸展是它半径的 \(\sqrt d\) 倍，而训练点的邻域并集在整个空间中的体积占比小到可以忽略。``在有限个点上正确''与``在这些点的邻域上正确''之间的距离，不会因为数据更多、模型更大而缩短。

保证只能来自分析，不能来自搜索。Lipschitz 认证与随机平滑给出的是不依赖攻击方法的半径，前者受限于常数难算、界太松，后者受限于推理时要采样、结论带置信水平。两者都保守，但两者说的话是算数的。

而工作点必须有人来选。标准准确率与鲁棒准确率之和有上界，这一点已经被构造证死；剩下的问题是在给定的应用里，哪个方向的错误更不能接受。这个问题没有数学答案。

上一个单元里，隐私与效用之间横着一条组合定理；这个单元里，准确与鲁棒之间横着一条构造性的不等式。两处的结构相同：想同时要两件好事，而数学给出了一条它们必须共同满足的约束，剩下的自由度只有选点。这一部分接下来的几个单元会在别的场景里重复这个形状——公平性的几个定义之间、在线决策的探索与利用之间、多方博弈的个体最优与集体最优之间，每一次都是同一件事换了一套符号。
